\chapter{Maxwell's Equations in Covariant Form}

Maxwell's four equations --- two homogeneous and two inhomogeneous --- govern the dynamics of the electromagnetic field itself. In this chapter we derive all four from the field tensor and show that they condense into just \emph{two} tensor equations, each elegant and manifestly Lorentz-covariant.

\section{From Potentials to Fields: a Brief Recap}

The electric and magnetic fields derive from the four-potential $A^\mu = (\phi, \vec{A})$ via the field tensor $F_{\mu\nu} = \partial_\mu A_\nu - \partial_\nu A_\mu$. In three-vector language:

\begin{align}
\vec{E} &= -\vec{\nabla}\phi - \frac{1}{c}\frac{\partial \vec{A}}{\partial t}, \label{eq:E_from_pot}\\
\vec{B} &= \vec{\nabla}\times\vec{A}. \label{eq:B_from_pot}
\end{align}

These definitions already encode two of Maxwell's equations, as we will now show.

\section{The Homogeneous Maxwell Equations}

\subsection{Derivation from the potential representation}

Consider the divergence of $\vec{B}$:

\begin{equation}
\vec{\nabla}\cdot\vec{B} = \vec{\nabla}\cdot(\vec{\nabla}\times\vec{A}) = 0,
\end{equation}

since the divergence of any curl vanishes identically. This is \textbf{Gauss's law for magnetism}: there are no magnetic monopoles.

Next, take the curl of $\vec{E}$:

\begin{equation}
\vec{\nabla}\times\vec{E} = -\vec{\nabla}\times\vec{\nabla}\phi - \frac{1}{c}\frac{\partial}{\partial t}(\vec{\nabla}\times\vec{A}).
\end{equation}

The first term vanishes (curl of a gradient is zero), and $\vec{\nabla}\times\vec{A} = \vec{B}$, giving

\begin{equation}
\vec{\nabla}\times\vec{E} = -\frac{1}{c}\frac{\partial\vec{B}}{\partial t}.
\end{equation}

This is \textbf{Faraday's law of induction}. Together, these two equations --- $\vec{\nabla}\cdot\vec{B} = 0$ and $\vec{\nabla}\times\vec{E} + \frac{1}{c}\frac{\partial\vec{B}}{\partial t} = 0$ --- are called the \emph{homogeneous} Maxwell equations because they have no source terms on the right-hand side.

\subsection{Covariant form: the Bianchi identity}

Both homogeneous equations are consequences of a single tensor identity. Define the \textbf{cyclic sum}

\begin{equation}
\partial_\lambda F_{\mu\nu} + \partial_\mu F_{\nu\lambda} + \partial_\nu F_{\lambda\mu}.
\end{equation}

Substituting $F_{\mu\nu} = \partial_\mu A_\nu - \partial_\nu A_\mu$:

\begin{align}
&\partial_\lambda(\partial_\mu A_\nu - \partial_\nu A_\mu) + \partial_\mu(\partial_\nu A_\lambda - \partial_\lambda A_\nu) + \partial_\nu(\partial_\lambda A_\mu - \partial_\mu A_\lambda) \nonumber\\
&= \partial_\lambda\partial_\mu A_\nu - \partial_\lambda\partial_\nu A_\mu + \partial_\mu\partial_\nu A_\lambda - \partial_\mu\partial_\lambda A_\nu + \partial_\nu\partial_\lambda A_\mu - \partial_\nu\partial_\mu A_\lambda.
\end{align}

Every term cancels with another (using the symmetry of mixed partial derivatives $\partial_\mu\partial_\nu = \partial_\nu\partial_\mu$), giving

\begin{equation}
\boxed{\partial_\lambda F_{\mu\nu} + \partial_\mu F_{\nu\lambda} + \partial_\nu F_{\lambda\mu} = 0.}
\end{equation}

This is the \textbf{Bianchi identity} for the electromagnetic field. It is an algebraic identity --- it holds for \emph{any} tensor of the form $F_{\mu\nu} = \partial_\mu A_\nu - \partial_\nu A_\mu$, regardless of the field equations.

\subsection{Equivalently, using the dual tensor}

The Bianchi identity can be written more compactly using the dual tensor $\widetilde{F}^{\mu\nu}$ introduced in Chapter~2:

\begin{equation}
\boxed{\partial_\mu \widetilde{F}^{\mu\nu} = 0.}
\end{equation}

This single equation, for each value of $\nu$, encodes both $\vec{\nabla}\cdot\vec{B} = 0$ (the $\nu = 0$ component) and Faraday's law (the $\nu = j$ components).

\subsection{Extracting the three-vector equations}

\textbf{Setting $\nu = 0$:}

\begin{equation}
\partial_\mu \widetilde{F}^{\mu 0} = \partial_i \widetilde{F}^{i0} = \partial_i B^i = \vec{\nabla}\cdot\vec{B} = 0,
\end{equation}

where we used $\widetilde{F}^{i0} = B_i$ from the explicit dual tensor matrix.

\textbf{Setting $\nu = j$:}

\begin{equation}
\partial_0 \widetilde{F}^{0j} + \partial_i \widetilde{F}^{ij} = 0.
\end{equation}

From the dual tensor, $\widetilde{F}^{0j} = -B_j$ and $\widetilde{F}^{ij} = \varepsilon^{ijk}E_k$. With $\partial_0 = \frac{1}{c}\frac{\partial}{\partial t}$:

\begin{equation}
-\frac{1}{c}\frac{\partial B_j}{\partial t} + (\vec{\nabla}\times\vec{E})_j = 0 \quad \Longrightarrow \quad \vec{\nabla}\times\vec{E} = -\frac{1}{c}\frac{\partial\vec{B}}{\partial t}.
\end{equation}

Both homogeneous Maxwell equations are recovered.

\section{The Inhomogeneous Maxwell Equations}

The homogeneous equations followed from the structure of $F_{\mu\nu}$ alone (the Bianchi identity). The \emph{inhomogeneous} equations are dynamical --- they relate the field to its sources: charge density $\rho$ and current density $\vec{J}$.

\subsection{The four-current}

In Gaussian units the \textbf{four-current density} is

\begin{equation}
J^\mu = (c\rho,\, \vec{J}),
\end{equation}

where $\rho$ is the charge density and $\vec{J}$ is the current density. The continuity equation (charge conservation) takes the covariant form

\begin{equation}
\partial_\mu J^\mu = \frac{\partial\rho}{\partial t} + \vec{\nabla}\cdot\vec{J} = 0.
\end{equation}

\subsection{The inhomogeneous field equation}

The two remaining Maxwell equations --- Gauss's law and the Amp\`ere--Maxwell law --- are contained in

\begin{equation}
\boxed{\partial_\mu F^{\mu\nu} = \frac{4\pi}{c}\,J^\nu.}
\end{equation}

This is the \textbf{inhomogeneous Maxwell equation} in covariant form (Gaussian units).

\subsection{Extracting the three-vector equations}

\textbf{Setting $\nu = 0$ (Gauss's law):}

\begin{equation}
\partial_\mu F^{\mu 0} = \partial_i F^{i0} = \partial_i E^i = \vec{\nabla}\cdot\vec{E} = \frac{4\pi}{c}\,J^0 = 4\pi\rho.
\end{equation}

Hence

\begin{equation}
\boxed{\vec{\nabla}\cdot\vec{E} = 4\pi\rho.}
\end{equation}

\textbf{Setting $\nu = j$ (Amp\`ere--Maxwell law):}

\begin{equation}
\partial_0 F^{0j} + \partial_i F^{ij} = \frac{4\pi}{c}\,J^j.
\end{equation}

From the contravariant tensor: $F^{0j} = -E_j$ and $F^{ij} = -\varepsilon^{ijk}B_k$. Then $\partial_0 F^{0j} = -\frac{1}{c}\frac{\partial E_j}{\partial t}$ and $\partial_i F^{ij} = (\vec{\nabla}\times\vec{B})_j$. Substituting:

\begin{equation}
-\frac{1}{c}\frac{\partial E_j}{\partial t} + (\vec{\nabla}\times\vec{B})_j = \frac{4\pi}{c}\,J_j,
\end{equation}

\begin{equation}
\boxed{\vec{\nabla}\times\vec{B} = \frac{4\pi}{c}\vec{J} + \frac{1}{c}\frac{\partial\vec{E}}{\partial t}.}
\end{equation}

The second term on the right is Maxwell's \emph{displacement current}, whose inclusion ensures consistency with charge conservation: taking the divergence of both sides and using $\vec{\nabla}\cdot(\vec{\nabla}\times\vec{B}) = 0$ and Gauss's law, one recovers $\partial_\mu J^\mu = 0$.

\section{Maxwell's Equations in Vacuum}

In regions with no charges or currents ($\rho = 0$, $\vec{J} = 0$, i.e.\ $J^\mu = 0$), Maxwell's equations in covariant form reduce to the beautifully symmetric pair:

\begin{align}
\partial_\mu F^{\mu\nu} &= 0, \label{eq:maxwell_vac_inhom}\\
\partial_\mu \widetilde{F}^{\mu\nu} &= 0. \label{eq:maxwell_vac_hom}
\end{align}

These two tensor equations replace the four three-vector equations. The first encodes $\vec{\nabla}\cdot\vec{E} = 0$ and $\vec{\nabla}\times\vec{B} = \frac{1}{c}\frac{\partial\vec{E}}{\partial t}$; the second encodes $\vec{\nabla}\cdot\vec{B} = 0$ and $\vec{\nabla}\times\vec{E} = -\frac{1}{c}\frac{\partial\vec{B}}{\partial t}$. In vacuum, the interchange $\vec{E} \to \vec{B}$, $\vec{B} \to -\vec{E}$ (electromagnetic duality) maps one equation into the other.

\section{Consistency: Charge Conservation from the Field Equations}

Taking the four-divergence of the inhomogeneous equation:

\begin{equation}
\partial_\nu\partial_\mu F^{\mu\nu} = \frac{4\pi}{c}\,\partial_\nu J^\nu.
\end{equation}

The left-hand side vanishes identically: $\partial_\nu\partial_\mu$ is symmetric in $\mu,\nu$ while $F^{\mu\nu}$ is antisymmetric, so their contraction is zero. Therefore

\begin{equation}
\partial_\nu J^\nu = 0,
\end{equation}

which is the continuity equation. Maxwell's equations are \emph{automatically consistent} with charge conservation --- this is not an additional assumption but a built-in consequence of the antisymmetry of $F^{\mu\nu}$.

\section{Summary: The Complete Set of Maxwell's Equations}

\begin{table}[h]
\centering
\renewcommand{\arraystretch}{1.6}
\begin{tabular}{lcc}
\toprule
\textbf{Equation} & \textbf{Covariant form} & \textbf{Three-vector form} \\
\midrule
Gauss (electric) & \multirow{2}{*}{$\partial_\mu F^{\mu\nu} = \dfrac{4\pi}{c}J^\nu$} & $\vec{\nabla}\cdot\vec{E} = 4\pi\rho$ \\
Amp\`ere--Maxwell & & $\vec{\nabla}\times\vec{B} = \dfrac{4\pi}{c}\vec{J} + \dfrac{1}{c}\dfrac{\partial\vec{E}}{\partial t}$ \\[6pt]
\midrule
Gauss (magnetic) & \multirow{2}{*}{$\partial_\mu\widetilde{F}^{\mu\nu} = 0$} & $\vec{\nabla}\cdot\vec{B} = 0$ \\
Faraday & & $\vec{\nabla}\times\vec{E} = -\dfrac{1}{c}\dfrac{\partial\vec{B}}{\partial t}$ \\
\bottomrule
\end{tabular}
\end{table}

Four vector equations have been compressed into two tensor equations. This is one of the great achievements of the covariant formulation: it makes the Lorentz invariance of electromagnetism manifest and reduces the bookkeeping from twelve component equations to eight (four per tensor equation, of which one is redundant by antisymmetry).
